\section{Optimization}
Test results for update processing (see chapter \ref{chap:results}) showed that
\snowman{} version 0.4 was very slow compared to the other software. The reason
for this was examined, and found to be a large amount of database calls spread
out in the code. Where \bhtb{} and \snortmanager{} used bulk insertion into the
database, \snowman{} 0.4 inserted objects as they were found and requested
objects as they were needed in the code. To solve this, the project group split
the methods of \pythoninline{Update} into two new classes: \pythoninline{Parser}
and \pythoninline{Updater}, where the first would only parse the incoming data
and send this to the latter which gathered all objects and wrote these in bulk
when the parsing was complete. As can be seen in chapter \ref{chap:results},
this had a huge impact on performance.

\begin{python}[caption=Updating a classification in \snowman{} 0.4]
from core.models import RuleClass

def updateClassification(type, description, priority):
  try:
    # Update existing classification
    ruleclass = RuleClass.objects.get(classtype=type)
    ruleclass.description = description
    ruleclass.priority = priority
    ruleclass.save()
  except RuleClass.DoesNotExist:
    # Add new classification
    ruleclass = RuleClass.objects.create(classtype=type, ...)
\end{python}